\documentclass[11pt,a4paper]{article}
\usepackage[margin=23mm,headheight=15pt]{geometry}
\usepackage{amsmath,amsthm,mathtools}
\usepackage{fontspec,unicode-math}
\setmainfont{lmroman10-regular.otf}[
  BoldFont=lmroman10-bold.otf,
  ItalicFont=lmroman10-italic.otf,
  BoldItalicFont=lmroman10-bolditalic.otf]
\setmathfont{latinmodern-math.otf}
\setmonofont{lmmono10-regular.otf}
\usepackage{booktabs,array,tabularx,enumitem,microtype,xcolor,graphicx}
\usepackage[unicode=true,colorlinks=true,linkcolor=blue!45!black,
  citecolor=blue!45!black,urlcolor=blue!45!black]{hyperref}
\usepackage{xurl,fancyhdr,needspace}
\hypersetup{
 pdftitle={Natural Core Motion: Face-Local Turning and Conditional Z Displacement},
 pdfauthor={Hilmir Frímann Halldórsson},
 pdfsubject={Research draft v0.1 - not yet author-approved or published}}
\pagestyle{fancy}
\fancyhf{}
\fancyhead[L]{\small Natural core motion: correspondence}
\fancyhead[R]{\small Research draft v0.1}
\fancyfoot[C]{\thepage}
\renewcommand{\headrulewidth}{0.3pt}
\setlength{\parindent}{0pt}
\setlength{\parskip}{4pt}
\setlength{\emergencystretch}{3em}
\setlist{itemsep=2pt,topsep=4pt,leftmargin=*}
\raggedbottom
\clubpenalty=10000
\widowpenalty=10000
\displaywidowpenalty=10000
\allowdisplaybreaks[2]
\numberwithin{equation}{section}
\newtheorem{proposition}{Proposition}[section]
\newtheorem{corollary}[proposition]{Corollary}
\theoremstyle{definition}
\newtheorem{definition}[proposition]{Definition}
\theoremstyle{remark}
\newtheorem{remark}[proposition]{Remark}
\newcommand{\R}{\mathbb R}
\newcommand{\C}{\mathbb C}
\newcommand{\iu}{\mathrm i}
\newcommand{\one}{\mathbf 1}
\newcommand{\eps}{\epsilon}
\newcommand{\Argz}{\operatorname{Arg}_0}
\newcommand{\re}{\operatorname{Re}}
\newcommand{\im}{\operatorname{Im}}
\newcommand{\diag}{\operatorname{diag}}
\newcommand{\norm}[1]{\left\lVert#1\right\rVert}
\newcommand{\abs}[1]{\left|#1\right|}
\newcommand{\Dnat}{\mathcal D_{\rm native}}
\newcommand{\Denv}{\mathcal D_{\rm envelope}}
\newcommand{\Dclock}{\mathcal D_{\rm clock}}
\newcommand{\envelope}{\mathcal E}
\newcommand{\muplace}{\mu_{\rm place}}

\begin{document}
\begin{center}
{\small\color{blue!45!black} TRI-OCTAGON MATHEMATICAL RESEARCH}\\[5mm]
{\LARGE\bfseries Natural Core Motion:\\[1mm]
Face-Local Turning and Conditional Z Displacement\par}
\vspace{5mm}
{\large Hilmir Frímann Halldórsson}\\[3mm]
29 October 2026\\[2mm]
{\small\bfseries Research draft v0.1 - not yet author-approved or published}
\end{center}

\begin{abstract}
The accepted Tri-Octagon recurrence evolves three complex channel values,
while its face decoder, axial attachment, and conditional Z placement assign
different geometric meanings to that state. We derive the motion consequences
of these existing maps without adding a material-core law. Matched transport
makes the decoded complex mean the collective vector in each face frame, and
its signed two-step turn transfers exactly to the corresponding outward face
normal. Common complex phase produces rotations within the three fixed face
planes; a pairwise-dot-product obstruction prevents a general interpretation
as one rigid ambient rotation. The staged Z placement is invariant under
that same common phase when observer settings are held fixed. An exact
telescoping identity separates native, envelope, and clock contributions to
its displacement. The completed coherent-domain theorem implies stationary
nonzero face vectors and a generally nonzero frozen-observer limit, whereas
the running staged envelope makes the placed point approach the center.
The accepted M3 cycle gives two alternating placed points at fixed observer
settings. These are deductions from inherited definitions and results. An
observer-only application design exposes the distinctions; the missing
relationship remains the assignment of native state information to the pose
of the central surface.
\end{abstract}

\textbf{Scientific status.}
This candidate develops exact algebraic consequences of accepted source
definitions and explicitly attributed prior theorems. It does not reopen
their certificates or publication dispositions. Software validation is
reported only when measured; Section~\ref{sec:implementation} records the
source implementation and its validation limits. No physical units, force identification, material
trajectory, or geometry-to-native feedback are adopted.

\section{Question, source scope, and distinct objects}\label{sec:question}

The current question is what motion within the accepted geometric
representation is already justified by the native equations. The recent
collective-turning investigation establishes a changing complex mean and,
on a separate domain, convergence to an unequal aligned profile
\cite{Collective}. Earlier work establishes the adopted face decoder and
matched connection \cite{Attachment}, the axial attachment and a certified
alternating cycle \cite{PaperG}, common-phase advance on a local chiral branch
\cite{Drift}, and a conditional placement of Z relative to the central
geometry \cite{Central}. Those results supply the starting maps.

Three objects must remain distinct. The collective native coordinate is
$m\in\C$. The refined central surface is a finite triangulated geometric
object $\Sigma_\tau$ in a fixed host, with a shape parameter $\tau$. A placed
observer is a point $X=\Phi(Z)\in\R^3$ obtained under a declared coordinate
placement. Neither the word ``core'' nor the common number of coordinates
identifies these objects with each other.

The scientific source snapshot used throughout is repository
\texttt{pzychozen/trioctagon-physics} at commit
\begin{center}\small
\texttt{7cc1fa03af589818bcdcfaaac51838cc36c84253}.
\end{center}
The references link to specific files at that commit. It is a reading
baseline, not a statement about a later working checkout. All simulations,
certificates, and publication results cited from earlier work retain their
original scopes. Our additions are the correspondence deductions and the
application contract below.

\section{Native state and the adopted geometric maps}\label{sec:definitions}

\subsection{Ordered native update}
Let $\Omega=x+\iu y\in\C^3$, with real channel triples $x,y$, and define
\begin{equation}
L_3=\one\one^T-3I,
\qquad m=\tfrac13\one^T\Omega,
\qquad \Pi=I-\tfrac13\one\one^T.
\end{equation}
For real native coefficients, the ordered map is
\begin{align}
V_i&=\Omega_i+\eps\Omega_i(k_i-|\Omega_i|^2)
                   +g(L_3\Omega)_i,\label{eq:pre}\\
\delta_i&=\lambda\sum_{j\ne i}
           \sin 3(\Argz V_j-\Argz V_i),\label{eq:kicks}\\
F_i(\Omega)&=V_i e^{\iu\delta_i}.\label{eq:post}
\end{align}
Every kick uses the same prestage $V$. The ideal convention is
$\Argz(0)=0$; synchronization is bypassed when $\lambda=0$. The native
update has no factor supplied by the observer's clock increment
\cite{Collective,Central}.

Common-phase covariance of $F$ is valid on the regular domain where the
relevant prestage entries are nonzero, and globally for the phase-off map.
It must not be extended through the native zero stratum without additional
analysis. In contrast, the pointwise decoder and readout identities proved
below apply to all finite channel states, including zeros.

\subsection{Fixed face frames and channel-area attachment}
The fixed outward unit normals and horizontal unit tangents are the columns
of the accepted matrices
\begin{equation}\label{eq:frames}
N=\begin{pmatrix}
-\sqrt3/2&0&\sqrt3/2\\
1/2&-1&1/2\\0&0&0
\end{pmatrix},\qquad
T=\begin{pmatrix}
-1/2&1&-1/2\\
-\sqrt3/2&0&\sqrt3/2\\0&0&0
\end{pmatrix}.
\end{equation}
With $e_z=(0,0,1)^T$, they satisfy
$t_i=e_z\times n_i$, $t_i\times e_z=n_i$, and
$t_i\cdot t_j=-1/2$ for distinct faces. The selected face centers $c_i$ are
fixed. The adopted decoder and matched transport are
\begin{equation}\label{eq:decoder}
D_i(z)=\re z\,t_i+\im z\,e_z\in W_i:=n_i^\perp,
\qquad T_{ij}=t_it_j^T+e_ze_z^T.
\end{equation}
The inverse on $W_i$ is
$E_i(v)=t_i\cdot v+\iu e_z\cdot v$. Thus $D$ retains the full native state.
Basing a vector at $c_i$ gives an attached arrow; its endpoint $c_i+D_i(z)$
need not belong to the finite polygon. The tangent plane is unbounded
\cite{Attachment}.

The signed channel area and its accepted attachment are
\begin{equation}\label{eq:chiral}
C=x\times y,\qquad W=TC,\qquad \Gamma=\one^TC.
\end{equation}
The vector $C$ lives in channel-index space. The specified $T$ gives $W$
its ambient planar axial-vector transformation law, while $\Gamma$ retains
the complementary common component. In particular,
\begin{equation}
T^TT=\tfrac32\Pi,\qquad
C=\tfrac23T^TW+\tfrac13\Gamma\one.
\end{equation}
The area convention is parallelogram area, with no extra factor $1/2$
\cite{PaperG}. These observables add no state variables or feedback.

\subsection{Staged Z and its conditional placement}
For the staged observer define
\begin{gather}
\kappa=\norm{\Omega},\qquad
K(\kappa)=a_{\rm vp}\frac{\kappa}{1+\kappa},\qquad
\envelope(t_c)=e^{-\gamma t_c},\label{eq:macrodefs}\\
d_\theta=(\cos\theta,\sin\theta,1)^T,\quad
f(\theta)=\cos3(\theta-\ell),\quad h_\theta=f(\theta)d_\theta,\notag\\
M=K(\kappa)\envelope(t_c)h_\theta,
\qquad Z=\alpha M+\beta C,\qquad
X=O+\muplace R_\psi Z.\label{eq:placement}
\end{gather}
Here $R_\psi$ is a fixed axial rotation and $\muplace>0$. The inherited
staged values are $a_{\rm vp}=0.618$, $\gamma=0.577$, $\ell=0.244$,
$\alpha=1$, and $\beta=1/2$, with
$\theta_n=(n\bmod12)\pi/6$ and $t_{c,n}=0.1n$. These decimal constants
remain the declared source parameters. The clock angle is distinct from the
native common phase. The historical EMA observer has a different scalar law
and is not substituted here \cite{Central,PaperE}.

The placement in \eqref{eq:placement} already supplies spatial points.
Its declaration does not make the raw channel action on $C$ an ambient
vector action under every shell symmetry; the accepted central analysis
computes the relevant residuals. Replacing $C$ by $W$ inside Z would define
a different observer \cite{Central}.

\begin{table}[htbp]
\centering\small
\caption{Correspondence and update dependence. No row introduces
geometry-to-native feedback.}
\begin{tabularx}{\linewidth}{@{}p{.19\linewidth}XXX@{}}
\toprule
Map and input & Output & Fixed data & What evolves\\\midrule
$D:\Omega\mapsto(v_i)$ & Three tangent vectors & Face frames and centers
  & Native state in an invertible representation\\
$T_{ij}:W_j\to W_i$ & Matched coefficients & Two face frames
  & Operator is fixed; no material carrier travels\\
$C\mapsto(W,\Gamma)$ & Axial plane vector and scalar & Matrix $T$
  & Passive native area readouts\\
$(\Omega,\theta,t_c)\mapsto X$ & Placed spatial point & $O,\muplace,R_\psi$
  & Native norm/area and observer schedule\\
$\tau\mapsto\Sigma_\tau$ & Refined finite surface & Declared host registration
  & No native update of $\tau$ or pose is defined\\
\bottomrule
\end{tabularx}
\label{tab:types}
\end{table}

\section{Collective turning transfers exactly to the face frames}
\label{sec:turn}

\begin{proposition}[Matched collective representation]\label{prop:mean}
For every finite state, if $v_i=D_i(\Omega_i)$, then
\begin{equation}\label{eq:matchedmean}
\boxed{D_i(m)=\frac13\sum_jT_{ij}v_j.}
\end{equation}
Moreover $T_{ij}T_{jk}=T_{ik}$ on $W_k$, and $D_i$ preserves the real
Euclidean norm and inner product of complex coefficients.
\end{proposition}
\begin{proof}
Orthonormality of $t_j,e_z$ gives $T_{ij}D_j(z)=D_i(z)$. Summing this
identity and using real linearity of $D_i$ proves \eqref{eq:matchedmean}.
Applying two transports retains the same two coefficients, proving the
composition identity. Finally
$D_i(z)\cdot D_i(w)=\re(\bar z w)$.
\end{proof}

Thus the mean is represented by transporting all face states into one
chosen frame and averaging. The raw ambient average is a different map:
\begin{equation}
\frac13\sum_iD_i(\Omega_i)
=\frac13Tx+(\im m)e_z.
\end{equation}
Since $T\one=0$, this ambient average discards the common real component.
It must not replace the matched average in a claim about $m$.

\begin{proposition}[Signed-turn transfer]\label{prop:turn}
For any $s,t\in\C$,
\begin{equation}\label{eq:turnidentity}
\boxed{D_i(s)\times D_i(t)=\im(\bar s t)\,n_i.}
\end{equation}
Consequently the signed two-increment turn of $m_n$ transfers exactly to
the matched collective vectors on every face, about that face's normal.
\end{proposition}
\begin{proof}
Write $s=a+\iu b$ and $t=c+\iu d$. The cross product is
$(ad-bc)t_i\times e_z=(ad-bc)n_i$, while $ad-bc=\im(\bar s t)$.
Apply the result to $s=m_{n+1}-m_n$ and $t=m_{n+2}-m_{n+1}$.
\end{proof}

The completed witness uses
\begin{equation}\label{eq:witness}
\eps=\frac1{10},\quad g=\frac1{60},\quad\lambda=\frac1{300},\quad
k=(0.999,1,1.001),\quad
\Omega_0=(1+\iu,1+\iu,1-2\iu).
\end{equation}
Its retained certificate gives
\begin{equation}
0.0110760731\le
\im\!\left[\overline{(m_1-m_0)}(m_2-m_1)\right]
\le0.0110760732.
\end{equation}
Proposition~\ref{prop:turn} therefore gives three corresponding positive
face-local turns. This application reuses the certified value; it neither
reruns the certificate nor recovers a historical trajectory
\cite{Collective}. The three normals are different axes in ambient space.

\begin{figure}[htbp]
\centering
\includegraphics[width=\linewidth]{figures/collective_turn.pdf}
\caption{The retained collective-turn witness and its exact face representation.
The numerical coordinates use midpoints of the published increment enclosures;
the inherited certificate supplies the sign. Face arrows use one fixed display
scale. Segments only connect sampled endpoints. No recurrence or certificate
was rerun to generate this illustration.}
\label{fig:turn}
\end{figure}

The statistic is oriented area, not angle alone. For nonzero increments it
equals the product of their lengths and the sine of their oriented angle.
If either increment is zero, the area remains defined but the angle is
undefined. Drawing lines between sampled vectors is a display convention;
the identity does not choose a continuous intervening trajectory.

\section{Local phase rotation and the rigid-congruence obstruction}
\label{sec:phase}

\begin{proposition}[Exact face-local phase action]\label{prop:phase}
For all $z\in\C$ and $\chi\in\R$,
\begin{equation}\label{eq:localphase}
\boxed{D_i(e^{\iu\chi}z)=
\cos\chi\,D_i(z)+\sin\chi\,n_i\times D_i(z).}
\end{equation}
The right side is the rotation of the tangent vector through $\chi$ about
the fixed normal $n_i$.
\end{proposition}
\begin{proof}
Since $n_i\times t_i=e_z$ and $n_i\times e_z=-t_i$, the right side has
coefficients $(\re z\cos\chi-\im z\sin\chi,
\re z\sin\chi+\im z\cos\chi)$ in the frame $(t_i,e_z)$. These are the
real and imaginary parts of $e^{\iu\chi}z$.
\end{proof}

In particular, the exact decoded native update is the already accepted
conjugate recurrence
\begin{align}
\widetilde v_i&=v_i+\eps(k_i-\norm{v_i}^2)v_i
       +g\sum_{j\ne i}(T_{ij}v_j-v_i),\label{eq:facepre}\\
v_i^+&=\cos\delta_i\,\widetilde v_i
       +\sin\delta_i\,n_i\times\widetilde v_i.\label{eq:facepost}
\end{align}
The phase kicks are those of the original prestage. This is $DFE$ on the
selected tangent-state space, not a second independently evolved model
\cite{Attachment}. It describes amplitude change, matched graph mixing,
and local rotation of attached vectors.

\begin{proposition}[No general rigid ambient phase action]
\label{prop:rigid}
The adopted decoder does not identify arbitrary common complex phase with
one rigid ambient rotation of the three decoded vectors. More strongly,
for any fixed state with at least two nonzero components, its full
continuous common-phase family is not a family of ambiently congruent
decoded vector triples.
\end{proposition}
\begin{proof}
A rigid ambient rotation preserves pairwise dot products. The state
$(1,1,1)$ decodes to $(t_1,t_2,t_3)$, whose distinct pairwise products are
$-1/2$. Multiplication by $\iu$ gives $(e_z,e_z,e_z)$, whose pairwise
products are $1$. This already disproves a general identification.

For the stronger statement, write $\Omega_i=r_ie^{\iu\phi_i}$ for each
nonzero component. For distinct $i,j$, elementary product-to-sum identities
give
\begin{align}
&D_i(e^{\iu\chi}\Omega_i)\cdot D_j(e^{\iu\chi}\Omega_j)\notag\\
&\quad =r_ir_j\left[
\frac14\cos(\phi_i-\phi_j)
-\frac34\cos(\phi_i+\phi_j+2\chi)\right].\label{eq:phasegram}
\end{align}
If $r_ir_j>0$, this is nonconstant as $\chi$ ranges continuously. No
orthogonal ambient transformation can change that pairwise product.
\end{proof}

This is a kinematic obstruction for the decoded vector configuration.
It is not an assertion that the native map takes the first example state
to the second, nor an exclusion of every special discrete rigid match.
It also does not prohibit a future, separately defined map from phase to
the pose of another geometric object. It establishes why such a map is not
already supplied by the face decoder.

\begin{figure}[htbp]
\centering
\includegraphics[width=\linewidth]{figures/phase_comparison.pdf}
\caption{A common-phase probe of the equal real state. The arrows rotate in
separate face planes; their pairwise dot products change, while the frozen Z
placement is invariant. The central surface is the fixed declared PLUS zero
member. This is a kinematic comparison, not a native trajectory.}
\label{fig:phase}
\end{figure}

\section{Common-phase advance and stationary placed observations}
\label{sec:invariance}

\begin{proposition}[Pointwise phase invariance]\label{prop:invariance}
For every finite state and every real $\chi$,
\begin{equation}
\kappa(e^{\iu\chi}\Omega)=\kappa(\Omega),\qquad
C(e^{\iu\chi}\Omega)=C(\Omega).
\end{equation}
Hence $W,\Gamma$, and the staged placement at fixed observer settings obey
\begin{equation}\label{eq:Xinvariant}
\boxed{X(e^{\iu\chi}\Omega;\theta,t_c)=X(\Omega;\theta,t_c).}
\end{equation}
\end{proposition}
\begin{proof}
The norm identity is immediate. Put $c=\cos\chi$ and $s=\sin\chi$.
The transformed real and imaginary triples are $x'=cx-sy$ and $y'=sx+cy$.
Their cross product is $(c^2+s^2)x\times y=C$. The remaining statements
follow by substitution into \eqref{eq:chiral} and
\eqref{eq:macrodefs}--\eqref{eq:placement}.
\end{proof}

The algebra includes zero components and does not use covariance of $F$.
This distinction is essential: a readout identity valid on the whole state
space does not repair a failure of native equivariance at a zero prestage.
Common-phase-invariant coherences and chiral observations on the regular
relative-equilibrium branch are already discussed in the drift source
\cite{Drift}. Equation~\eqref{eq:Xinvariant} is their direct consequence
under the existing staged placement.

On the accepted regular D1 branch,
$F(\Omega)=e^{\iu\omega}\Omega$. The decoded face vectors therefore
rotate locally, while their norm, chirality and frozen-observer Z remain
constant. If the observer schedule advances, its Z points can still change
through the envelope and clock. D1 concerns a specified local chiral branch
with a saddle reference; it supplies neither attraction of arbitrary inputs
nor a material-core rotation law \cite{Drift}.

The face representation and the placed observer retain different
information. A constant placed point does not imply a stationary native
state. Conversely, a nonzero value of an axial observable is not already an
angular velocity or displacement per update. The full native state remains
the object that is stepped.

\section{Exact attribution of placed displacement}\label{sec:attribution}

Let $K_n=K(\kappa_n)$, $E_n=\envelope(t_{c,n})$, and
$h_n=h_{\theta_n}$. Keep $O,\muplace,R_\psi,\alpha,\beta$ fixed during a
comparison.

\begin{proposition}[Native, envelope, and clock telescope]
\label{prop:telescope}
For finite states with $\Omega_{n+1}=F(\Omega_n)$,
\begin{equation}\label{eq:telescope}
\boxed{X_{n+1}-X_n=
\muplace R_\psi(\Dnat+\Denv+\Dclock),}
\end{equation}
where
\begin{align}
\Dnat&=\alpha E_nh_n(K_{n+1}-K_n)
                    +\beta(C_{n+1}-C_n),\label{eq:nativepart}\\
\Denv&=\alpha K_{n+1}(E_{n+1}-E_n)h_n,\label{eq:envelopepart}\\
\Dclock&=\alpha K_{n+1}E_{n+1}(h_{n+1}-h_n).
\label{eq:clockpart}
\end{align}
The native term is precisely
\begin{equation}\label{eq:nativedifference}
\muplace R_\psi\Dnat
=X(F\Omega_n;\theta_n,t_{c,n})
 -X(\Omega_n;\theta_n,t_{c,n}).
\end{equation}
\end{proposition}
\begin{proof}
Subtract the two instances of \eqref{eq:placement}, and insert intermediate
macro products $K_{n+1}E_nh_n$ and $K_{n+1}E_{n+1}h_n$. The three
differences are \eqref{eq:nativepart}--\eqref{eq:clockpart}. Holding the
observer settings fixed gives \eqref{eq:nativedifference} directly.
\end{proof}

This is the placed, general-weight form of the existing staged attribution
\cite[Section 8.2]{Central}. Its order is explicit: first change the native
state, then the envelope, then the angle. Another telescoping order can
allocate mixed finite changes differently, while leaving the total
unchanged. Thus the stated native contribution is a well-defined
fixed-observer comparison; the three labels do not claim an order-free
physical force decomposition.

For pure common-phase advance, Proposition~\ref{prop:invariance} makes
$\Dnat=0$. For a general orbit, both its norm and channel area contribute.
Those contributions can cancel. The value of $\Dnat$ alone does not
reconstruct the full state change. If placement or display parameters are
changed between samples, \eqref{eq:telescope} no longer accounts for the
whole displayed difference unless that additional change is included
separately or the trace is restarted.

\section{Consequences of the completed convergence and cycle results}
\label{sec:consequences}

\subsection{The coherent domain}
The completed collective paper proves the following inherited result at
the coefficients in \eqref{eq:witness}. The domain with
$0.9\le|\Omega_i|\le1.1$ and all ordinary phase lifts in one interval of
width at most $\pi/12$ is invariant. Every orbit there converges to
\begin{equation}\label{eq:inheritedlimit}
\Omega_n\longrightarrow e^{\iu\theta_\infty}r_*.
\end{equation}
Here $r_*$ is the unique positive aligned fixed amplitude profile in
$[0.9,1.1]^3$, its entries are unequal, and $\theta_\infty$ depends on the
initial state. Moreover, if $\Delta_0$ is the initial ordinary phase width,
\begin{equation}\label{eq:Cbound}
|C_{i,n}|\le1.1^2\Delta_0 c^n,
\qquad c=1-\frac{3\sqrt2}{200}<1.
\end{equation}
These are reused theorem hypotheses and bounds, not new certificates
\cite{Collective}. The turning entrance of \eqref{eq:witness} is outside
that domain; no capture connection is assumed.

\begin{corollary}[Face and placed-observer limits]\label{cor:limits}
Under the preceding coherent-domain hypotheses,
\begin{equation}\label{eq:facelimit}
D_i(\Omega_{i,n})\longrightarrow
r_{*,i}(\cos\theta_\infty\,t_i+\sin\theta_\infty\,e_z).
\end{equation}
Thus all decoded vectors settle to fixed nonzero values, while
$C_n,W_n,\Gamma_n\to0$.
For fixed observer settings,
\begin{equation}\label{eq:frozenlimit}
X_n\longrightarrow
O+\muplace R_\psi
 \left[\alpha K(\norm{r_*})\envelope(t_c)h_\theta\right].
\end{equation}
If instead $\gamma>0$ and $t_{c,n}\to+\infty$, the staged placement
satisfies $X_n\to O$, regardless of whether its angle sequence converges.
\end{corollary}
\begin{proof}
Continuity and real linearity of $D_i$ give \eqref{eq:facelimit}.
At the aligned limit, the real and imaginary channel triples are scalar
multiples of $r_*$, so their cross product is zero. Equation
\eqref{eq:Cbound} gives the explicit decay, and the fixed attachments are
continuous. Substitution gives \eqref{eq:frozenlimit}.

For the running envelope, $K(\kappa_n)\le a_{\rm vp}$ and
$\norm{h_{\theta_n}}\le\sqrt2$. Since $R_\psi$ is orthogonal,
\begin{equation}\label{eq:placedbound}
\norm{X_n-O}\le\muplace\left[
 |\alpha|\sqrt2\,a_{\rm vp}e^{-\gamma t_{c,n}}
 +|\beta|\sqrt3\,(1.1)^2\Delta_0c^n\right].
\end{equation}
Both terms tend to zero.
\end{proof}

The frozen-observer limit is generally a nonzero macro offset. Every
possible limiting common phase gives the same point in
\eqref{eq:frozenlimit}, although it gives different tilts in
\eqref{eq:facelimit}. Under the staged schedule, approach of $X_n$ to the
center combines native area decay with the imposed macro envelope. It
does not show that the central surface translates inward. The positive
residual spread $\norm{\Pi r_*}>0$ is also compatible with zero area:
unequal collinear complex coefficients have those two properties.

\subsection{The M3 limiting two-cycle}
Paper G separately proves capture, under its own specified coefficients
and entrance, by a full-state cycle with
\begin{equation}\label{eq:M3}
\Omega^-=F(\Omega^+)=e^{\iu\vartheta}\overline{\Omega^+},
\qquad F(\Omega^-)=\Omega^+.
\end{equation}
Its limiting $W$ and $\Gamma$ are nonzero and change sign at each step.
The approaching orbit converges by parity; it is not an exact two-cycle at
every earlier sample. The local onset theorem and this finite capture
theorem retain their separate scopes \cite{PaperG}.

\begin{corollary}[Two placed snapshots at a fixed observer]
\label{cor:M3}
On the exact cycle \eqref{eq:M3}, let $C^+=C(\Omega^+)$ and
$\kappa_*=\norm{\Omega^+}$. With observer settings fixed, put
\begin{equation}
X_{\rm macro}=O+\muplace R_\psi
     [\alpha K(\kappa_*)\envelope(t_c)h_\theta].
\end{equation}
Then
\begin{equation}\label{eq:twopoints}
\boxed{X^\pm=X_{\rm macro}\pm\muplace\beta R_\psi C^+,}
\qquad
\norm{X^+-X^-}=2\muplace|\beta|\norm{C^+}.
\end{equation}
\end{corollary}
\begin{proof}
Conjugation preserves the norm and reverses $C$; a subsequent common phase
preserves both norm and $C$. Hence the two macro values agree and the
chiral terms are opposite. Orthogonality of $R_\psi$ gives the distance.
\end{proof}

Equation~\eqref{eq:twopoints} supplies two alternating points and their
separation. It supplies no continuous half-turn, gap-crossing path, or
accumulated spatial rotation. An advancing observer can change the macro
offset, so its scheduled total need not itself be a two-cycle. Likewise,
the reversal $W\mapsto-W$ is a signed observable reversal, not a measured
angular velocity.

\section{Central geometry and the missing pose correspondence}
\label{sec:pose}

The accepted folded host, refined family $\Sigma_\tau$, and preserved
historical intersecting mesh are different constructions. The central
analysis declares the refined finite faces and their registration; this
paper uses that geometry as context and does not replace its connectivity
\cite{Central}. In particular, $\tau$ is a shape parameter, not the native
update index, observer time, or common phase.

An inherited contact constraint is already informative. An axial rotation
and translation preserving all three exact upper contacts of the registered
surface has zero translation and angle in $2\pi\mathbb Z/3$. A continuous
motion in that class is therefore constant. This is a conditional geometric
constraint, not an assertion that contacts are welded. The proper half-turn
symmetry at the zero family member is an endpoint operation and does not
provide a continuous registered path \cite[Section 5]{Central}.

The central fan openings and finite host-notch patches are distinct sets.
Matching azimuth, meeting an infinite plane, membership in a finite
triangle, and passage through an opening are different claims. Neither the
discrete native update nor the affine placement chooses an interpolation
or a transport law through those openings. Contact claims must continue to
use the accepted finite-region tests.

The smallest remaining relationship is the response of central orientation
to common native phase. The decoder retains that phase as a local directional
pattern, whereas the staged Z placement discards it at fixed observer
settings. If the central surface is intended to turn under pure common-phase
advance, a phase-sensitive pose assignment must be supplied and checked
against its geometry and contacts. If its pose is intended to remain
unchanged, amplitude and relative-state changes are the relevant inputs.
The present mathematics distinguishes these choices but does not select
one on the author's behalf.

A passive pose assignment would itself still differ from material dynamics
and from feedback into $F$. No such assignment, physical force, or change
to the native recurrence is made in this work. The established result is
the exact face-motion correspondence together with conditional observer
displacement.

\section{Observer-only application design and verification scope}
\label{sec:implementation}

The application feature specified by this paper is a motion-correspondence
view of the existing recurrence. Its canonical evolving value is $\Omega$,
stepped through the existing native implementation. The geometry and
readouts are derived from that state. The feature is intended to show:
\begin{enumerate}
\item decoded channel arrows and optional matched collective arrows on the
fixed face frames;
\item the sampled complex mean and its signed-turn diagnostic;
\item conditional placed Z samples, including a fixed-observer comparison;
\item the three contributions in \eqref{eq:telescope}; and
\item a separately labeled common-phase probe applied to a display copy.
\end{enumerate}
The central surface serves as fixed geometric reference. The phase probe
demonstrates \eqref{eq:localphase} and \eqref{eq:Xinvariant}; it is not a
native trajectory and must not silently modify the stored evolving state.
An implementation can expose meaningful geometric motion without adding
an unsupported animation of a solid core.

The plotted samples must retain the distinction between logical vectors
and display scaling. A fixed common arrow scale multiplies the displayed
turn area by its square; reported scientific statistics should use the
unscaled values. Per-sample normalization can change apparent paths and
does not preserve the turn calculation. At a zero vector, direction is
undefined. At a zero increment, its inter-increment angle is undefined.
These cases should be reported explicitly rather than assigned a fictitious
orientation.

\begin{table}[htbp]
\centering\small
\caption{Bounded verification protocol. These are required comparisons,
not a report of executed tests.}
\begin{tabularx}{\linewidth}{@{}p{.31\linewidth}X@{}}
\toprule
Comparison & Expected mathematical property\\\midrule
Actual face frames & $t_i\times e_z=n_i$ and the accepted Gram products\\
Collective arrows & Agreement of decoded mean and matched transported average\\
Signed turn & Agreement of complex area and each decoded cross product\\
Phase probe & Moving local arrows with invariant $\kappa,C,W,\Gamma,X$ at fixed observer\\
Rigid-congruence example & Pairwise dot products distinguish the two triples\\
Displacement attribution & Native, envelope and clock contributions sum to the placed difference\\
State ownership & Probe and observer operations leave native state and later native stepping unchanged\\
Degenerate readouts & Undefined directions and angles remain distinct from numerical zero\\
\bottomrule
\end{tabularx}
\end{table}

\subsection{Implemented component and measured validation}
The source candidate adds \texttt{motion\_observer.py},
\texttt{motion\_panel.py}, and \texttt{motion\_view.py} to the existing
scientific application, with a small integration in \texttt{app.py}. The
new panel is in workspace B, Observers and Diagnostics. It computes only
after explicit analysis of a selected, already validated three-channel
Core RunRecord. The cached result is a separate application observation;
it is not inserted into the Core record history. Scrubbing selects cached
rows. The kernel, worker protocol, distribution locks, and prior papers
are not modified.

The component preserves the actual stored update indices, including gaps.
Its interval quantities are endpoint differences over the reported update
span. For a sparse interval, the native counterfactual replaces the old
state by the later stored state at fixed observer settings; it need not be
one application of $F$. Its display observer is recomputed at every stored
state, including the first. This is explicitly separate from historical
constructor-zero rows and from all recorded staged or EMA outputs.

The view uses the hash-checked existing folded-host reference and the
fixed PLUS $\tau=0$ central surface. One arrow scale is chosen for the full
analyzed history. Placement is fixed within a completed observation;
changing its scale or rotation requires explicit reanalysis. A common-phase
probe transforms a copy of the whole cached history. Source changes and
pending placement edits disable export; the old cache remains identifiable.
Exports use the application's existing Windows file-handle publication
mechanism and do not overwrite occupied destinations.

The independent numerical suite has 23 passing tests, and the offscreen
panel suite has 13 passing tests, including synchronous rendering of all
three plot views with and without a phase probe. These counts denote
software checks, not independent mathematical theorems. The algebra tests
include a hand-calculated signed-turn example, midpoint consistency with
the retained witness enclosure, phase-invariant zero-component states,
independently evaluated ordered counterfactual positions, sparse indices,
degenerate directions, the 18-vertex/42-edge/24-triangle fixed surface,
immutable caches, source preservation, input refusal, and export lineage.
The panel tests cover explicit calculation, cached-only navigation,
source-context guards, pending settings, failed operations, exports, and
accessible numerical alternatives.

On the tested phase examples, the largest absolute residual was
$8.89\times10^{-16}$ for $\Gamma$, and at most $1.12\times10^{-16}$
for the placed points. Attribution reconstruction had residual at most
$1.12\times10^{-16}$; comparison with separately evaluated ordered
counterfactual positions had residual at most $4.45\times10^{-16}$.
These rounded-up descriptions summarize measured binary64 residuals;
they are not interval certificates or universal error bounds.

The executed environment is Linux x86-64, Python 3.12.14, and
PySide6 Essentials 6.11.2 with offscreen Qt. Plot checks used Matplotlib
3.10.8; the unchanged Windows application pin is 3.11.2 and remains to be
tested in its supported lane. A main-window integration
smoke check used the archived application fixture and a test-only lock-loader
stub; it preserved the dynamics draft and submitted no scientific job.
The supported Windows/Python 3.11 installed lane and real Windows
file-handle export behavior have not been executed in this environment.
The delivered candidate includes targeted Windows validation instructions.
This source integration is not a certified wheel, an installed application
replacement, or a publication release.

\section{Conclusions and retained programme directions}

The existing maps provide several mathematically distinct kinds of motion.
The complex collective turn is exactly a matched local turn in every face
frame. Phase is a local rotation operation on those vectors, with a direct
obstruction to a general rigid ambient interpretation. The staged placed
observer is insensitive to common phase, and its finite displacement has
an exact native/envelope/clock decomposition. Completed convergence and
cycle theorems consequently produce different geometric readout behavior:
stationary nonzero face vectors, conditional approach of Z to the center,
or two alternating placed snapshots under a frozen observer.

These results support a faithful observational application feature and
locate the next mathematical choice precisely: the map from native state
information to the central surface's pose. They do not yet select that map.
The broader gravity--magnetism question, central Z transport, the proposed
outer geometric structure, legitimate SRG integration, and the separate
TORMENT and embodiment programme remain standing directions. Three-Way
remains parked. No result here integrates those directions automatically.

\Needspace{7\baselineskip}
\paragraph{AI-assisted preparation disclosure.}
This research candidate was prepared through human--AI collaboration,
including AI-assisted mathematical exposition, source reconciliation and
LaTeX drafting. The author supplies the original geometric direction and
retains scientific design and publication judgment. This document is not
yet author-approved or published, and no independent peer review is claimed.

\begin{thebibliography}{9}
\small
\bibitem{Attachment}
H. F. Halldórsson, \emph{Native Response Geometry: equivariant spatial
observables, state--shell attachment, and closure}, accepted source chapters,
2026. Pinned source:
\url{https://github.com/pzychozen/trioctagon-physics/blob/7cc1fa03af589818bcdcfaaac51838cc36c84253/papers/NATIVE_RESPONSE_GEOMETRY/manuscript/observables_closure.tex}.

\bibitem{Collective}
H. F. Halldórsson, \emph{Transient Collective Turning and Coherent
Relaxation in the Native Triad}, v0.1, 2026. Published source:
\url{https://github.com/pzychozen/trioctagon-physics/blob/7cc1fa03af589818bcdcfaaac51838cc36c84253/papers/COLLECTIVE_TURNING_AND_COHERENT_RELAXATION/v0.1/main.tex}.

\bibitem{Central}
H. F. Halldórsson, \emph{Central Geometry, Z Readouts, and Chiral Response
in the Tri-Octagon Model}, publication v0.3, 2026. Pinned source:
\url{https://github.com/pzychozen/trioctagon-physics/blob/7cc1fa03af589818bcdcfaaac51838cc36c84253/papers/CENTRAL_GEOMETRY_Z_CHIRAL_RESPONSE/publication/v0.3/manuscript.tex}.

\bibitem{PaperG}
H. F. Halldórsson, \emph{Geometry-Attached Axial Observables,
Period-Doubling, and Certified Entrance Capture in the Tri-Octagon Map},
Paper G v0.1.1, 2026. Pinned source:
\url{https://github.com/pzychozen/trioctagon-physics/blob/7cc1fa03af589818bcdcfaaac51838cc36c84253/papers/PAPER_G/v0.1.1/source/paper_g.tex}.

\bibitem{Drift}
H. F. Halldórsson, \emph{Cubic Common-Phase Drift on a Chiral Branch of
the Tri-Octagon Map}, research publication v0.1, 2026. Pinned source:
\url{https://github.com/pzychozen/trioctagon-physics/blob/7cc1fa03af589818bcdcfaaac51838cc36c84253/papers/MEASUREMENT_AND_CHIRAL_DRIFT/v0.1/drift/manuscript.tex}.

\bibitem{PaperE}
H. F. Halldórsson, \emph{Paper E: the historical Z manifold and its
distinct observer definitions}, v0.1.1, 2026. Pinned source:
\url{https://github.com/pzychozen/trioctagon-physics/blob/7cc1fa03af589818bcdcfaaac51838cc36c84253/papers/PAPER_E/v0.1.1/PAPER_E_Z_MANIFOLD_v0.1.1.md}.
\end{thebibliography}
\end{document}
